\begin{tuluyen}
Trong một nhà máy sản xuất amonia theo quy trình Haber-Bosch, giai đoạn sản xuất khí hydrogen bằng phản ứng của methane và nước được thực hiện theo phương trình hóa học (1) sau:
\[
\ce{CH4(g) + H2O(g) ->[xt] CO(g) + 3H2(g)} \tag{1}
\]
Phản ứng (1) là phản ứng thu nhiệt mạnh. Lượng nhiệt này được cung cấp từ quá trình đốt cháy hoàn toàn khí methane theo phương trình hóa học (2):
\[
\ce{CH4(g) + 2O2(g) ->[t^\circ] CO2(g) + 2H2O(g)} \tag{2}
\]
Xét các phản ứng ở điều kiện chuẩn và hiệu suất chuyển hóa của methane là 100\%. Tính khối lượng khí methane cần thiết để có thể sản xuất được 3 tấn \ce{H2} trong giai đoạn trên. Biết rằng: ở phản ứng (1) để phản ứng hoàn toàn 1 mol \ce{CH4} cần cung cấp lượng nhiệt bằng 205,8 kJ; ở phản ứng (2) khi đốt cháy hoàn toàn 1 mol \ce{CH4} tỏa ra lượng nhiệt bằng 802,5 kJ, trạng thái của chất có kí hiệu (g) -- chất khí. Giả sử 90\% lượng nhiệt tỏa ra từ phản ứng (2) được cung cấp cho phản ứng (1).
\source{Trích nguyên văn Câu 4.3, đề tuyển sinh chuyên Hóa Điện Biên 2026.}
\dapso{\(\approx\SI{10.28}{\ton}\) methane.}
\begin{loigiaiBT}
Số mol hydrogen cần sản xuất:
\[
n_{\ce{H2}}=\frac{3\times10^6}{2}=\SI{1.5e6}{\mole}.
\]
Theo phản ứng (1), số mol methane dùng để tạo hydrogen:
\[
n_{\ce{CH4},1}=\frac{1.5\times10^6}{3}=\SI{5.0e5}{\mole}.
\]
Nhiệt lượng phản ứng (1) cần:
\[
Q_{\text{cần}}=5.0\times10^5\times205.8=\SI{1.029e8}{\kilo\joule}.
\]
Gọi $x$ là số mol methane cần đốt ở phản ứng (2). Vì chỉ $90\%$ nhiệt lượng được cung cấp cho phản ứng (1):
\[
x\times802.5\times0.90=1.029\times10^8.
\]
Suy ra:
\[
x\approx\SI{1.42471e5}{\mole}.
\]
Tổng số mol methane:
\[
n_{\text{tổng}}\approx5.0\times10^5+1.42471\times10^5
=\SI{6.42471e5}{\mole}.
\]
Khối lượng methane:
\[
m=6.42471\times10^5\times16\approx\SI{1.02795e7}{\gram}
\approx\SI{10.28}{\ton}.
\]
\textbf{Vậy,} cần khoảng \SI{10.28}{\ton} methane.
\end{loigiaiBT}
\end{tuluyen}
